\begin{solucion}{Comparacion de contratos}
Definimos los escenarios:
\[
\begin{array}{c|l}
\toprule
\text{Escenario} & \text{Significado} \\
\midrule
D & \text{sin contrato, cadena descentralizada.}\\
A & \text{contrato de reparto con } \alpha.\\
B & \text{contrato con tarifa fija } F.\\
\bottomrule
\end{array}
\]
Para cada escenario \(e\in\{D,A,B\}\):
\[
\begin{array}{c|l}
\toprule
\text{Simbolo} & \text{Significado} \\
\midrule
\Pi_s^e & \text{utilidad del proveedor en el escenario } e.\\
\Pi_r^e & \text{utilidad del minorista en el escenario } e.\\
\Pi_{SC}^e & \text{utilidad total de la cadena: } \Pi_s^e+\Pi_r^e.\\
\bottomrule
\end{array}
\]
Ademas:
\[
\begin{array}{c|l}
\toprule
\text{Parametro} & \text{Significado} \\
\midrule
\alpha\in(0,1) & \text{fraccion asignada al minorista en el contrato A.}\\
F\in\mathbb{R}_+ & \text{pago fijo del minorista al proveedor en el contrato B.}\\
\bottomrule
\end{array}
\]

Primero recordamos los resultados sin contrato:
\[
\Pi_s^D=1225.125,
\qquad
\Pi_r^D=612.5625,
\qquad
\Pi_{SC}^D=1837.6875.
\]

La utilidad centralizada maxima es:
\[
\Pi_{SC}^C=2450.25.
\]

\textbf{Contrato A: reparto de utilidades.}

Con \(\alpha=0.35\):
\[
w=\alpha c=0.35.
\]
El contrato coordina, por lo que la utilidad total es:
\[
\Pi_{SC}^{A}=2450.25.
\]
La utilidad del minorista:
\[
\Pi_r^{A}=\alpha\Pi_{SC}^{A}
=0.35(2450.25)
=857.5875.
\]
La utilidad del proveedor:
\[
\Pi_s^{A}=(1-\alpha)\Pi_{SC}^{A}
=0.65(2450.25)
=1592.6625.
\]

\textbf{Contrato B: tarifa fija.}

Si \(w=c=1\), el minorista enfrenta el costo marginal real de la cadena. Por lo tanto, elige la cantidad centralizada:
\[
q_C=49.5,
\qquad
p_C=50.5,
\qquad
\Pi_{SC}^{B}=2450.25.
\]
Antes de pagar la tarifa fija, toda la utilidad operacional queda en el minorista:
\[
2450.25.
\]
Como paga \(F=1400\), queda:
\[
\Pi_r^B=2450.25-1400=1050.25.
\]
El proveedor recibe:
\[
\Pi_s^B=F=1400.
\]

\textbf{Resumen:}
\[
\begin{array}{l|ccc}
\toprule
\text{Escenario} & \Pi_s & \Pi_r & \Pi_{SC} \\
\midrule
\text{Sin contrato} & 1225.13 & 612.56 & 1837.69 \\
\text{Contrato A: } \alpha=0.35 & 1592.66 & 857.59 & 2450.25 \\
\text{Contrato B: } F=1400 & 1400.00 & 1050.25 & 2450.25 \\
\bottomrule
\end{array}
\]

Ambos contratos A y B coordinan la cadena, porque ambos logran:
\[
\Pi_{SC}=2450.25.
\]
Pero distribuyen la utilidad de forma distinta. El proveedor prefiere A:
\[
1592.66>1400.
\]
El minorista prefiere B:
\[
1050.25>857.59.
\]

La respuesta tecnica es: ambos contratos son eficientes para la cadena, pero la eleccion depende del poder de negociacion de las partes.
\end{solucion}
